\cotefile{85}
\scanpage{2}
% L1
\uline{Question} du centre de $\Gamma=G_2/{-}$

\noindent\textcircled{1}\(\left\{\vphantom{\rule{0pt}{30mm}}\right.\)\begin{minipage}{.88\linewidth}
% L2
Représentations \inserted{n. dég.} \deleted{de} $G_0=\mathbb Z_{/2}$

% L3
Représentations \inserted{simplem\unsure{ent}} dégénérées de $\mathbb D_\infty=G_1$ -- i.e. représ. \inserted{fidèles} de $\mathbb D_2$ / relation
% L4
avec oct-aèdre, ou quaternions / quaternions généralisés\ldots.

% L5
Représentations simplement dégénérées de $G_2$ -- i.e.
% L6
représentations n.\ dég. de $\mathbb D_\infty\times\mathbb Z_{/2}$ [\deleted{considérer}\inserted{déterminer}
% L7
les cartes correspondantes, et leur genre\ldots.]

% L8
Représentations doublem\unsure{ent} dégénérées de $G_2$ -- i.e. repr.
% L9
de $\mathbb D_2\times\mathbb Z_{/2}\simeq{\mathbb Z_{/2}}^3$ \deleted{avec} de ``rang'' \deleted{\unsure{égal}} $\geqslant 2$
\end{minipage}

\noindent\rule{.33\linewidth}{.4pt}

% L10
\textcircled{\scriptsize $1'$}\quad$\Vert\!\mid$\quad Pts de ramification des repr. \deleted{de} $G_2$ \ldots

\noindent\rule{.33\linewidth}{.4pt}

% margin note written sideways
\margin{cf verso}
\noindent\(\left\{\vphantom{\rule{0pt}{34mm}}\right.\)\begin{minipage}{.91\linewidth}
% L11
Groupes et cartes associés aux représentations \inserted{projectives}
% L12
de $G_2$ \deleted{dans} un corps -- notamment sur
% L13
un corps fini

\noindent\rule{.60\linewidth}{.4pt}

% L14
Les rev. \deleted{de degré}\inserted{étales modérés} de genre $0$ de la
% L15
sphère privée de trois points

\noindent\rule{.60\linewidth}{.4pt}

% L16
Les cartes finies associées aux hyperpolygônes
% L17
réguliers finis (cf verso)
\end{minipage}

\noindent\rule{.37\linewidth}{.4pt}

\noindent\textcircled{2}\(\left\{\vphantom{\rule{0pt}{16mm}}\right.\)\begin{minipage}{.87\linewidth}
% L18
Terminer de regarder les syst. de \unsure{Schwartz|Schwarz} pythagoriciens
% L19
(cas de $G_5$), et déterminer
% L20
\deleted{\unsure{l'explicitation|l'explication}} la structure de leurs : op. sous les
% L21
gpes des op. élémentaires, dans les trois cas
\end{minipage}

\noindent\rule{.30\linewidth}{.4pt}

\noindent\textcircled{4}\(\left\{\vphantom{\rule{0pt}{10mm}}\right.\)\begin{minipage}{.87\linewidth}
% L22
structure des gpes des op. élémentaires,
% L23
et relations avec $\mathrm{Gal}(\overline{\mathbb Q}/\mathbb Q)$
\end{minipage}

\noindent\rule{.41\linewidth}{.4pt}

% L24
\textcircled{3}\quad Détermination des hyperpolygônes pythagoriciens \unsure{(leurs} \deleted{\unsure{auto}} \unsure{d'équations|l'équation|l'ambiguïté}\unsure{$\{\ldots\}$)}
% archive sheet number: 1
